% Part: second-order-logic
% Chapter: syntax-and-semantics
% Section: semantic-notions

\documentclass[../../../include/open-logic-section]{subfiles}

\begin{document}

\olfileid{sol}{syn}{sem}
\olsection{المفاهيم الدلالية}

\begin{explain}
تعرف المفاهيم المنطقية المركزية، وهي \emph{الصحة} و\emph{الاستلزام}
و\emph{قابلية الإرضاء}، في منطق الرتبة الثانية بالطريقة نفسها التي
تعرف بها في منطق الرتبة الأولى، غير أن علاقة الإرضاء الأساسية صارت
الآن علاقة إرضاء~!!{formula}s الرتبة الثانية. وبالطبع تكون
!!{sentence} الرتبة الثانية~!!a{formula} ربطت فيها جميع المتغيرات،
ومنها المتغيرات المحمولية والدالية.
\end{explain}

\begin{defn}[الصحة]
تكون الجملة~$!A$ \emph{صحيحة منطقيًا}، $\Entails !A$، إذا وفقط إذا
$\Sat{M}{!A}$ لكل~!!{structure}~$\Struct M$.
\end{defn}

\begin{defn}[الاستلزام]
\emph{تستلزم} مجموعة الجمل~$\Gamma$ الجملة~$!A$،
$\Gamma \Entails !A$، إذا وفقط إذا، لكل~!!{structure}~$\Struct M$
تحقق~$\Sat{M}{\Gamma}$، تحقق أيضًا~$\Sat{M}{!A}$.
\end{defn}

\begin{defn}[قابلية الإرضاء]
تكون مجموعة الجمل~$\Gamma$ \emph{قابلة للإرضاء} إذا تحقق
$\Sat{M}{\Gamma}$ لبعض~!!{structure}~$\Struct M$. وإذا لم تكن
$\Gamma$ قابلة للإرضاء سميت \emph{غير قابلة للإرضاء}.
\end{defn}

\end{document}
